\chapter{Literature Review}

\section{Introduction}
This chapter reviews the work that the thesis builds on and defines the components
it modifies. The literature is grouped by theme rather than listed paper by paper:
first the progression of object detectors, then the techniques developed
specifically for small and aerial targets, and then human detection in disaster
imagery. A separate section explains the base architectures (YOLO11, the CBAM
and ECA attention modules, multi-scale detection, and state-space models) so
that Chapter~III can concentrate on the proposed modifications. The chapter ends
with the research gap that the thesis addresses.

\section{Literature Review}

\subsection{From Two-Stage Detectors to Real-Time and Transformer Models}
Modern object detection divides into two-stage and one-stage families. Two-stage
detectors, from R-CNN~\cite{girshick2014rcnn} to Faster R-CNN~\cite{ren2015fasterrcnn}
and its Cascade variant~\cite{cai2018cascade}, first propose regions and then
classify them, which gives strong accuracy at a high computational cost. One-stage
detectors such as SSD~\cite{liu2016ssd} predict boxes directly and trade a little
accuracy for speed; RetinaNet~\cite{lin2017focal} addressed the resulting
class-imbalance with the focal loss, and the YOLO family, from the original
design~\cite{redmon2016yolo} through successive versions~\cite{redmon2018yolov3,%
bochkovskiy2020yolov4,wang2023yolov7,jocher2023yolov8,wang2024yolov9,wang2024yolov10}
to YOLO11~\cite{jocher2024yolo11}, made single-pass detection the practical standard
for real-time use. A third, transformer-based family removes hand-designed
components such as anchor boxes and non-maximum suppression: DETR~\cite{carion2020detr}
posed detection as set prediction, and RT-DETR~\cite{zhao2024rtdetr} later made this
competitive with YOLO at real-time speeds. This thesis works within the one-stage
YOLO family because its speed and small model size suit airborne deployment.

\subsection{Small-Object and Aerial Detection}
Detecting very small objects is a distinct problem, surveyed in detail by Cheng et
al.~\cite{cheng2023smallobject}, because targets occupy few pixels and are lost in
the coarse strides of a standard detection head; drone benchmarks such as
VisDrone~\cite{zhu2021visdrone} make this the dominant difficulty. Three lines of
work address it. The first adds higher-resolution detection scales: feature
pyramids~\cite{lin2017fpn} and path-aggregation necks fuse fine and coarse features,
and adding a stride-4 (P2) head extends detection to targets a few pixels wide;
TPH-YOLOv5~\cite{zhu2021tphyolov5} combined such a head with transformer prediction
for drone imagery. The second re-weights features with attention so that faint
targets are not suppressed by background clutter; channel attention
(ECA~\cite{wang2020eca}) and combined channel-and-spatial attention
(CBAM~\cite{woo2018cbam}) are the common choices. The third operates at inference
time: slicing-aided hyper inference (SAHI~\cite{akyon2022sahi}) runs the detector
over overlapping crops so that small targets are effectively enlarged. All three are
relevant here; the thesis studies the first two as architectural changes and the
third as an inference-time setting.

\subsection{Label Assignment for Tiny Objects}
A fourth line of work changes neither the network nor the input but the training
step that decides which predictions count as positive samples. Its starting
observation is that IoU-based matching breaks down for tiny boxes: a one-pixel
shift between a candidate and a ground-truth box a few pixels wide can drive
their IoU to zero, so the assigner ranks candidates for the smallest targets
almost arbitrarily. Wang et al.~\cite{wang2021nwd} replaced the IoU with a
normalised Wasserstein distance between Gaussians fitted to the boxes, which
decays smoothly with offset instead of collapsing, and reported the largest
gains when the substitution is made in the assignment step rather than in the
regression loss. RFLA~\cite{xu2022rfla} built a full assignment strategy on a
receptive-field Gaussian prior, and SimD~\cite{shi2024simd} learned an adaptive
location-and-shape similarity for the same purpose. On the tiny-object
benchmarks these methods target, the assignment-side change is consistently
what moves sub-16-pixel accuracy; loss-only variants underperform. Modern YOLO
versions select their training samples with the task-aligned assigner of
TOOD~\cite{feng2021tood}, which still scores localisation with the IoU family,
and none of the tiny-object assignment work above has been applied inside it
for aerial person detection. That gap is the opening for the assignment method
of Chapter~III.

\subsection{Human Detection in Disaster and Search-and-Rescue Imagery}
Aerial human detection for search and rescue has been surveyed as a benchmark
problem in its own right~\cite{sar_survey}, and prior work has applied deep models
to aerial video for this purpose~\cite{aerial_human_video}. A recurring obstacle is
data: real annotated disaster imagery with very small human targets is scarce. Real
search-and-rescue datasets such as SARD~\cite{sambolek2021sard} and
HERIDAL~\cite{bozic2019heridal} are valuable but limited in scale and scene variety.
Nihal et al.~\cite{nihal2024c2a} addressed this with C2A, a semi-synthetic dataset
that composites human poses, segmented with a U$^2$-Net model~\cite{qin2020u2net},
onto disaster backgrounds drawn from aerial emergency
imagery~\cite{kyrkou2020aider}, and benchmarked a range of detectors on it. Their
study is the direct point of comparison for this thesis, which uses the same dataset
and test split.

Training on composited or synthetic data raises the question of transfer to real
imagery. Work on domain randomisation showed that models trained on varied
synthetic scenes can transfer to the real world~\cite{tobin2017domainrand}, and
that fine-tuning a synthetic-trained model on a small amount of real data
outperforms training on the real data alone~\cite{tremblay2018synthetic}. The
practical recipe that follows, and the one examined in this thesis, is to train
on the large synthetic set and then adapt with whatever real labels can be
afforded, measuring what remains of the gap afterwards.

\subsection{State-Space Models for Vision}
State-space models have recently been proposed as a linear-time alternative to
attention for long-sequence modelling. Building on structured state
spaces~\cite{gu2022s4}, Mamba~\cite{gu2023mamba} introduced a selective scan that
makes the state dynamics input-dependent. Vision adaptations followed: Vision
Mamba~\cite{zhu2024visionmamba} and VMamba~\cite{liu2024vmamba} scan image patches in
multiple directions, and Mamba-YOLO~\cite{wang2025mambayolo} applied the idea to
detection by replacing the backbone with state-space blocks. The variant explored in
this thesis is narrower: it keeps
the convolutional backbone and inserts bidirectional local-window state-space
blocks only in the neck, which lets its effect be measured against an otherwise
identical model.

\section{Core Concepts}
This section describes the base architectures that the thesis modifies. The
proposed changes to them are the subject of Chapter~III.

\subsection{The YOLO11 Baseline}
YOLO11~\cite{jocher2024yolo11} is a single-stage detector with three parts: a
convolutional backbone that extracts features at decreasing spatial resolution
through C3k2 stages, a path-aggregation neck (PAN--FPN) that fuses features across
scales, and a decoupled detection head that predicts boxes at three strides (P3,
P4, P5). The YOLO11 backbone ends with a spatial-pyramid-pooling (SPPF) stage
followed by a C2PSA self-attention block. The medium variant, YOLO11m, is used as
the baseline throughout this thesis.

\subsection{Attention: CBAM and ECA}
Attention modules re-weight feature maps so that informative channels and spatial
locations are emphasised. Squeeze-and-Excitation networks~\cite{hu2018senet} first
popularised channel re-weighting; ECA~\cite{wang2020eca} makes it lightweight with a
one-dimensional convolution. CBAM~\cite{woo2018cbam} is stronger:
it applies channel attention followed by spatial attention in sequence, as shown in
Figure~\ref{fig:cbam}. An earlier study on this dataset compared both against the
baseline and found CBAM to give better operational recall than ECA at a similar
cost, which is why CBAM is used in this thesis; that comparison is summarised in
Chapter~IV.

\begin{figure}[tbp]
  \centering
  \figorplaceholder{figures/fig_cbam_module.png}{6.5cm}
  \caption{The CBAM module. Channel attention (from global average and max pooling
  through a shared multi-layer perceptron) refines the feature map, after which
  spatial attention (from channel-wise pooling and a $7\times7$ convolution)
  refines it again.}
  \label{fig:cbam}
\end{figure}

\subsection{Multi-Scale Detection and the P2 Scale}
The finest standard detection scale, P3 at stride~8, summarises an $8\times8$ pixel
region per grid cell at $640$ input, which is coarse for targets only a few pixels
across. Adding a P2 scale at stride~4 introduces a $160\times160$ feature map whose
cells cover $4\times4$ pixels, extending detection to targets as small as about
four pixels. Figure~\ref{fig:p2} shows how the P2 branch attaches alongside the
standard P3--P5 scales.

\begin{figure}[tbp]
  \centering
  \figorplaceholder{figures/fig_p2_head.png}{6.5cm}
  \caption{The four-scale detection head. The P2 scale (stride~4, $160\times160$)
  is added to the standard P3--P5 scales to resolve very small targets; each scale
  feeds per-scale detections into non-maximum suppression.}
  \label{fig:p2}
\end{figure}

\subsection{State-Space Models and the C3k2Mamba Block}
A state-space model maps an input sequence to an output through a hidden state
whose dynamics, in the selective (Mamba) formulation, depend on the input. Applied
to vision, a feature map is scanned as a sequence. The block explored in this
thesis, C3k2Mamba (Figure~\ref{fig:mamba}), scans a local window in both the forward
and reverse directions and fuses the two, and is inserted in place of selected C3k2
neck blocks. Its effect on this task is measured in Chapter~IV.

\begin{figure}[tbp]
  \centering
  \figorplaceholder{figures/fig_c3k2mamba.png}{6.5cm}
  \caption{The explored C3k2Mamba block: forward and reverse selective scans over a
  local window are fused and projected, replacing the C3k2 bottleneck at selected
  neck layers while leaving the backbone and head unchanged.}
  \label{fig:mamba}
\end{figure}

\section{Discussion}

\subsection{Research Gap and Proposed Solution}
The reviewed work leaves a specific gap. Small-object techniques, namely extra
detection scales, attention, and inference-time slicing, are usually reported in
isolation and on generic aerial data, and the recent enthusiasm for state-space
necks has not been tested against a controlled baseline for very small human
targets. On disaster imagery specifically, the C2A benchmark reports whole-model
scores but not a component-wise account of where a small-object gain actually comes
from. This thesis fills that gap with a single-variable additive ablation on C2A
that measures attention, the P2 detection scale, and a state-space neck against one
another under an identical training and evaluation protocol, and then reports which
components are worth their cost. The four configurations are summarised in
Figure~\ref{fig:chain}, and the specific differences from prior work are listed in
Table~\ref{tab:gap}. Two further gaps are addressed in the extended study. The
tiny-object assignment line has not been carried into the task-aligned assigner
that modern YOLO versions actually use, and published disaster-detection work
rarely measures its models on real disaster footage at all, let alone on an
event held out entirely from training; Chapter~IV does both.

\begin{figure}[tbp]
  \centering
  \figorplaceholder{figures/fig_ablation_chain.png}{4.5cm}
  \caption{The additive ablation studied in this thesis: each step changes one
  component of the YOLO11m baseline, isolating the effect of attention, the P2
  detection scale, and the state-space neck.}
  \label{fig:chain}
\end{figure}

\begin{table}[tbp]
  \centering
  \caption{Position of this thesis relative to prior work on small-object aerial
  human detection.}
  \label{tab:gap}
  \begin{tabular}{@{}p{4.2cm} p{5cm} p{5cm}@{}}
    \toprule
    Aspect & Prior work & This thesis \\
    \midrule
    Small-object techniques & reported individually & compared in one controlled ablation \\
    Attention choice & asserted or single option & CBAM vs ECA measured on the task \\
    State-space neck & proposed as a backbone, on generic data & tested in the neck against an identical baseline; reported honestly \\
    Target domain & mostly generic aerial imagery & disaster/SAR imagery (C2A), very small human targets \\
    Tiny-object assignment & shown on tiny-object benchmarks, outside YOLO's assigner & blended into the task-aligned assigner; matched pair with confidence intervals \\
    Real-footage validation & rarely reported for disaster models & own drone, campus, and disaster sets, with a held-out unseen event \\
    \bottomrule
  \end{tabular}
\end{table}
